\chapter{Literature Survey}

\section{Introduction}

The literature survey is an essential phase of software development that involves studying existing systems, research contributions, and technologies related to the proposed project. It helps in understanding current approaches, identifying their strengths and limitations, and recognizing the research gap that motivates the development of a new solution. For the PrepSmart platform, the literature survey focuses on modern placement preparation platforms, online learning systems, coding practice portals, interview preparation resources, and educational technologies that support campus recruitment.

Over the past decade, digital learning platforms have significantly transformed placement preparation by providing students with online coding environments, structured learning materials, mock assessments, interview experiences, and skill development resources. These platforms have improved accessibility to quality educational content and enabled students to prepare independently from any location. However, most existing systems focus on specific aspects of placement preparation rather than providing a unified learning ecosystem.

Students preparing for campus recruitment typically rely on multiple platforms for coding practice, aptitude preparation, interview guidance, resume development, and progress tracking. This fragmented approach often results in inconsistent preparation, inefficient resource management, and difficulty in evaluating overall placement readiness. Furthermore, many platforms provide limited personalization and lack integrated analytical features that help students identify strengths and weaknesses throughout their preparation journey.

The objective of this chapter is to examine widely used placement preparation platforms, analyze their features, identify their limitations, and establish the necessity for a comprehensive platform such as PrepSmart. The findings of this survey provide the foundation for the design and development decisions presented in the subsequent chapters.

\section{Review of Existing Placement Preparation Platforms}

The increasing demand for technical skills has resulted in the emergence of several online platforms that support placement preparation. These platforms provide resources such as programming challenges, interview experiences, technical articles, aptitude questions, and online courses. Although they have significantly improved accessibility to educational content, most of them focus on specific aspects of placement preparation rather than providing a comprehensive solution. The following subsections present a review of widely used placement preparation platforms.

\subsection{LeetCode}

LeetCode is one of the most widely used online coding platforms for technical interview preparation. It provides a large collection of programming problems covering data structures, algorithms, databases, system design, and competitive programming. The platform also offers company-specific interview questions, coding contests, and discussion forums that help students prepare for software engineering interviews.

Despite its extensive collection of coding problems, LeetCode primarily focuses on programming practice and interview questions. It does not provide aptitude preparation, resume management, personalized learning paths, interview simulations, or integrated placement analytics, requiring students to depend on additional platforms for complete preparation.

\subsection{HackerRank}

HackerRank is an online platform designed for programming practice, skill assessment, and technical recruitment. It supports multiple programming languages and provides challenges related to algorithms, databases, artificial intelligence, and software development. Organizations also use HackerRank to conduct online coding assessments during recruitment.

Although HackerRank offers an excellent environment for coding practice and technical assessments, it provides limited support for interview preparation, learning progress analysis, resume management, and company-specific preparation. Students must therefore combine HackerRank with other learning resources to prepare comprehensively.

\subsection{GeeksforGeeks}

GeeksforGeeks is a popular educational platform that provides tutorials, interview preparation articles, coding questions, aptitude resources, and explanations of computer science concepts. It serves as a valuable reference for students preparing for technical interviews by offering extensive theoretical content and programming tutorials.

However, GeeksforGeeks primarily functions as a learning resource rather than an integrated placement preparation platform. Features such as personalized dashboards, structured preparation tracking, AI-assisted interview practice, and unified progress monitoring are limited.

\subsection{InterviewBit}

InterviewBit focuses on technical interview preparation through guided learning paths, coding exercises, and interview-oriented programming questions. The platform emphasizes structured learning and provides company-specific interview experiences that assist students in understanding real recruitment processes.

While InterviewBit offers organized learning content and coding practice, it lacks several complementary features such as resume management, portfolio development, AI-assisted interviews, integrated analytics, and comprehensive placement tracking.

\subsection{LinkedIn Learning}

LinkedIn Learning is an online educational platform offering professional courses across technology, business, communication, and leadership domains. Students can access expert-led video courses to strengthen both technical and soft skills required for career development.

Although LinkedIn Learning provides high-quality educational content, it is not specifically designed for placement preparation. It does not include integrated coding practice, mock assessments, company-specific preparation, or placement performance analytics within a unified platform.

\section{Comparative Analysis of Existing Platforms}

Table \ref{tab:comparison} presents a comparative analysis of widely used placement preparation platforms based on important features required during campus recruitment. The comparison highlights that while existing platforms perform well in specific domains such as coding practice or technical learning, none of them provide a unified ecosystem that integrates learning, assessments, coding practice, interview preparation, resume management, and analytical performance tracking within a single application.

\begin{table}[H]
\centering
\caption{Comparison of Existing Placement Preparation Platforms}
\label{tab:comparison}
\renewcommand{\arraystretch}{1.25}
\resizebox{\textwidth}{!}{
\begin{tabular}{|p{5cm}|c|c|c|c|c|c|}
\hline
\textbf{Feature} &
\textbf{LeetCode} &
\textbf{HackerRank} &
\textbf{GFG} &
\textbf{InterviewBit} &
\textbf{LinkedIn Learning} &
\textbf{PrepSmart} \\
\hline

Coding Practice & \cmark & \cmark & \cmark & \cmark & \xmark & \cmark \\ \hline

Structured Learning Paths & \xmark & Limited & \cmark & \cmark & \cmark & \cmark \\ \hline

Aptitude Preparation & \xmark & Limited & \cmark & Limited & \xmark & \cmark \\ \hline

Computer Science Subjects & Limited & Limited & \cmark & \cmark & \cmark & \cmark \\ \hline

Mock Tests & \xmark & \cmark & Limited & \cmark & \xmark & \cmark \\ \hline

AI Interview Preparation & \xmark & \xmark & \xmark & \xmark & \xmark & \cmark \\ \hline

Resume Management & \xmark & \xmark & \xmark & \xmark & \xmark & \cmark \\ \hline

Portfolio Management & \xmark & \xmark & \xmark & \xmark & \xmark & \cmark \\ \hline

Company-Specific Preparation & Limited & Limited & Limited & \cmark & \xmark & \cmark \\ \hline

Performance Analytics & Limited & Limited & Limited & Limited & \xmark & \cmark \\ \hline

Progress Dashboard & \xmark & Limited & \xmark & Limited & \xmark & \cmark \\ \hline

Single Integrated Platform & \xmark & \xmark & \xmark & \xmark & \xmark & \cmark \\ \hline

\end{tabular}
}
\end{table}

\section{Research Gap}

The comparative analysis demonstrates that existing placement preparation platforms primarily focus on individual aspects of learning rather than providing a comprehensive preparation environment. Coding platforms such as LeetCode and HackerRank emphasize programming practice, while educational platforms such as GeeksforGeeks and LinkedIn Learning concentrate on theoretical learning and professional development. InterviewBit offers structured interview preparation but lacks several complementary features required for complete placement readiness.

A significant limitation observed across these platforms is the absence of an integrated ecosystem that combines structured learning, coding practice, aptitude preparation, interview preparation, resume management, portfolio development, company-specific preparation, and analytical dashboards within a single application. Students are therefore required to use multiple independent platforms, resulting in fragmented learning experiences, inefficient progress tracking, and inconsistent preparation strategies.

Another important research gap is the limited use of intelligent analytics for monitoring placement readiness. Existing systems generally provide isolated performance metrics without offering comprehensive insights into learning progress, topic-wise strengths and weaknesses, or overall employability. Similarly, AI-assisted interview preparation and unified career management features remain largely unavailable across existing platforms.

The identified research gap establishes the need for a comprehensive and intelligent placement preparation platform capable of integrating multiple preparation activities into a single ecosystem. PrepSmart has been developed to address these limitations by providing a unified platform that combines learning resources, coding practice, assessments, interview preparation, resume and portfolio management, company-specific preparation, and personalized performance analytics.

\section{Chapter Summary}

This chapter presented a detailed review of existing placement preparation platforms and analyzed their strengths and limitations. A comparative study highlighted that although platforms such as LeetCode, HackerRank, GeeksforGeeks, InterviewBit, and LinkedIn Learning provide valuable learning resources, they primarily focus on specific aspects of placement preparation rather than offering a comprehensive solution.

The analysis identified the absence of a unified platform capable of integrating structured learning, coding practice, assessments, AI-assisted interview preparation, resume management, portfolio development, company-specific preparation, and performance analytics within a single ecosystem. These observations establish the motivation for developing PrepSmart as an integrated and scalable placement preparation platform.

The next chapter presents the system analysis and design of the proposed solution, including requirement analysis, software architecture, database design, UML diagrams, and overall system workflow.